\documentclass[11pt,a4paper]{article}
% Spolecny styl pro ucitelske pripravy MAT (Work/MAT)
% Verze 1.1, 2026-10-23
\NeedsTeXFormat{LaTeX2e}
\ProvidesPackage{priprava}[2026/10/23 v1.1 Sablona ucitelskych priprav MAT]

% Zakladni typografie a cestina
\RequirePackage[T1]{fontenc}
\RequirePackage[utf8]{inputenc}
\RequirePackage[czech]{babel}
\RequirePackage{lmodern}
\RequirePackage{microtype}
\RequirePackage{amsmath,amssymb}
\RequirePackage{xcolor}
\RequirePackage{enumitem}
\RequirePackage[most]{tcolorbox}
\RequirePackage{tabularx,booktabs,array}
\RequirePackage{fancyhdr}
\RequirePackage{needspace}
\RequirePackage{geometry}
\RequirePackage[hidelinks]{hyperref}

\geometry{a4paper,top=17mm,bottom=18mm,left=23mm,right=23mm}
\definecolor{notegray}{gray}{0.39}
\setlength{\parindent}{0pt}
\setlength{\parskip}{2.5pt plus 1pt minus 0.5pt}
\setlist[itemize]{topsep=2pt,itemsep=1pt,parsep=0pt,leftmargin=1.45em}
\setlist[enumerate]{topsep=2pt,itemsep=2pt,parsep=0pt,leftmargin=1.65em}
\setcounter{secnumdepth}{0}
\renewcommand{\arraystretch}{1.12}

% Zahlavi: vlevo rocnik | globalni tema, vpravo cislo: nazev hodiny
% Zapati zustava zcela prazdne.
\pagestyle{fancy}
\fancyhf{}
\renewcommand{\headrulewidth}{0.35pt}
\renewcommand{\footrulewidth}{0pt}
\setlength{\headheight}{14pt}
\newcommand{\nastavPripravu}[4]{%
  \fancyhead[L]{\small #1\enspace\textbar\enspace #2}%
  \fancyhead[R]{\small #3: #4}%
}

% Minimalisticke cernobile boxy. Kazda definice zvlast.
\tcbset{commonlesson/.style={enhanced,breakable=false,colback=white,colframe=black,
  colbacktitle=white,coltitle=black,boxrule=0.45pt,
  arc=0.8mm,outer arc=0.8mm,left=3mm,right=3mm,top=1.6mm,bottom=1.6mm,
  toptitle=1mm,bottomtitle=0.7mm,fonttitle=\bfseries,
  attach boxed title to top left={xshift=2.5mm,yshift=-1.15mm},
  boxed title style={colback=white,colframe=white,boxrule=0pt,left=1mm,right=1mm}}}
\newtcolorbox{zadaniBox}[1][]{commonlesson,title={Zadání},#1}
\newtcolorbox{prikladBox}[1][]{commonlesson,title={Příklad},#1}
\newtcolorbox{definiceBox}[1][]{commonlesson,title={Definice},#1}

% Spolecny box pro zadani a jeho reseni, bez mezery mezi castmi.
% Pouziti:
%   \begin{prikladReseniBox}
%     Zadani ulohy.
%     \tcblower
%     \begin{align*} ... \end{align*}
%   \end{prikladReseniBox}
% Spodni cast ma jemny sedy podklad a delici linku; lze v ni sazet
% vice radku matematiky. Obe casti maji jen jeden spolecny ramecek.
\tcbset{solutionlesson/.style={
  commonlesson,
  enhanced jigsaw,
  breakable,
  colbacklower=black!3,
  segmentation style={solid,draw=black!25,line width=0.35pt},
  fontlower=\small\color{notegray},
  before lower={\noindent\textbf{Řešení:}\enspace},
  leftlower=3mm,rightlower=3mm
}}
\newtcolorbox{prikladReseniBox}[1][]{solutionlesson,title={Příklad},#1}
\newtcolorbox{zadaniReseniBox}[1][]{solutionlesson,title={Zadání},#1}

% Pomucky pro autora pripravy: sedive poznamky; \reseni pro starsi pripravy
\newcommand{\poznamka}[1]{\par\begingroup\small\color{notegray}\itshape
  \textnormal{\textbf{Poznámka pro učitele.}}\enspace #1\par\endgroup}
\newcommand{\reseni}[1]{\par\begingroup\small\color{notegray}
  \textbf{Očekávané řešení:}\enspace #1\par\endgroup}
\newcommand{\blok}[1]{\needspace{5\baselineskip}\subsection*{#1}}
\newcommand{\vyrok}[1]{\uv{#1}}

\nastavPripravu{1. ročník}{Algebraické výrazy}{47}{Dělení mnohočlenu jednočlenem}
\begin{document}

% Zdroje: vlastni ucebnice Algebraicke vyrazy, Deleni mnohoclenu jednoclenem;
% Realisticky.cz: 010805 Dělení mnohočlenů (metodická inspirace).
\poznamka{Jádro 45 minut: zpětná vazba na násobení 5 min, pravidlo a podmínky 10 min, dva vzorové příklady 10 min, samostatná práce 15 min, kontrola 5 min. Rezerva je volitelná.}
\blok{Dělení jako obrácené násobení}
Z rovnosti $2x(x^2+3x+1)=2x^3+6x^2+2x$ můžeme zpětně poznat, čím se každý člen násobil. Dělení součtu jednočlenem převede celý výraz na součet podílů.
\begin{zadaniBox}[title={Pravidlo dělení mnohočlenu jednočlenem}]
Je-li $M\ne0$ jednočlen a $P=A_1+\cdots+A_k$ mnohočlen, pak všude, kde je podíl definován, platí
\[\frac{P}{M}=\frac{A_1}{M}+\cdots+\frac{A_k}{M}.\]
Každý člen dělíme zvlášť. Při dělení mocnin stejného základu exponenty odečítáme.
\end{zadaniBox}
\begin{prikladReseniBox}[breakable=false,title={Příklad 1: podíl je opět mnohočlen}]
Uprav $\dfrac{6x^4-9x^3+3x^2}{3x^2}$ a uveď podmínky.
\tcblower
Pro $x\ne0$:
\[\frac{6x^4}{3x^2}-\frac{9x^3}{3x^2}+\frac{3x^2}{3x^2}=2x^2-3x+1.\]
Zkouška: $(3x^2)(2x^2-3x+1)=6x^4-9x^3+3x^2$.
\end{prikladReseniBox}
\begin{prikladReseniBox}[breakable=false,title={Příklad 2: výsledkem nemusí být mnohočlen}]
Uprav $\dfrac{9x^3-18x^2+27x-\frac19}{-9x}$.
\tcblower
Pro $x\ne0$ platí
\[-x^2+2x-3+\frac{1}{81x}.\]
Výsledný výraz není mnohočlen, protože obsahuje proměnnou ve jmenovateli.
\end{prikladReseniBox}
\poznamka{Zkrácený výsledek může mít širší definiční obor než původní zadání; podmínku $x\ne0$ proto nesmažeme. Příkladem je $x^2/x=x$ pouze pro $x\ne0$.}
\blok{Samostatné procvičení — základ}
\begin{zadaniReseniBox}[breakable=false,title={Samostatná práce: vyděl a urči přípustné hodnoty}]
\begin{enumerate}[label=\alph*)]
\item $\dfrac{8x^3+4x^2}{4x}$;
\item $\dfrac{-12a^4+6a^2}{-3a^2}$;
\item $\dfrac{6x^3y^2-9x^2y}{3xy}$;
\item $\dfrac{4x^2+2x+1}{2x}$;
\item $\dfrac{x^3-x^2}{x^2}$.
\end{enumerate}
\tcblower
\begin{enumerate}[label=\alph*)]
\item $2x^2+x$, $x\ne0$;
\item $4a^2-2$, $a\ne0$;
\item $2x^2y-3x$, $x\ne0,\ y\ne0$;
\item $2x+1+\frac{1}{2x}$, $x\ne0$;
\item $x-1$, $x\ne0$.
\end{enumerate}
\end{zadaniReseniBox}
\blok{Volitelná rezerva}
\begin{zadaniReseniBox}[breakable=false,title={Rozšíření: najdi chybu v krácení}]
Žák uvádí $\dfrac{x^2+1}{x}=x+1$. Rozhodni, zda má pravdu.
\tcblower
Ne. Pro $x\ne0$ je správně $x+\frac1x$. Každý člen čitatele se dělí zvlášť; navíc $x=0$ je vyloučeno.
\end{zadaniReseniBox}
\blok{Domácí procvičení}
\begin{zadaniReseniBox}[breakable=false,title={Domácí úkol}]
\begin{enumerate}[label=\alph*)]
\item $\dfrac{15x^4-10x^3+5x^2}{5x^2}$;
\item $\dfrac{-8a^3b^2+12a^2b}{4ab}$;
\item $\dfrac{6x^2+3}{3x}$;
\item $\dfrac{2x^3-4x^2}{-2x}$;
\item $\dfrac{3x^2y-6xy^2+9xy}{3xy}$;
\item $\dfrac{5x^3+10x}{5x}$.
\end{enumerate}
\tcblower
\begin{enumerate}[label=\alph*)]
\item $3x^2-2x+1$, $x\ne0$;
\item $-2a^2b+3a$, $a,b\ne0$;
\item $2x+\frac1x$, $x\ne0$;
\item $-x^2+2x$, $x\ne0$;
\item $x-2y+3$, $x\ne0$, $y\ne0$;
\item $x^2+2$, $x\ne0$.
\end{enumerate}
\end{zadaniReseniBox}
\end{document}
